\input{../tex-inputs/latex-leseansicht-vorspann}

\section[Arthur Schnitzler und Richard Beer-Hofmann an Gustav Schwarzkopf, 9. 8. 1899]{L04435 Arthur Schnitzler und Richard Beer-Hofmann an Gustav Schwarzkopf, 9. 8. 1899}
\nopagebreak\mylabel{L04435v}
\rehead{ }\normalsize\beginnumbering\briefempfaengerindex{Schwarzkopf, Gustav@\textsc{Schwarzkopf, Gustav}!zzzBeer-Hofmann, Richard@\emph{von Richard Beer-Hofmann}!1899-08-091@{9. 8. 1899}|(be}\briefempfaengerindex{Schwarzkopf, Gustav@\textsc{Schwarzkopf, Gustav}!zzzSchnitzler, Arthur@\emph{von Arthur Schnitzler}!1899-08-091@{9. 8. 1899}|(be}
\toendnotes[C]{\smallbreak\pagebreak[2]}
\correspDesc{Versand  durch Arthur Schnitzler, Richard Beer-Hofmann am 9. 8. 1899 in Predazzo
\newline{}Übermittlung  am 9. 8. 1899 in Predazzo
\newline{}Zustellung  am 11. 8. 1899 in Wien 1/1 1
\newline{}Erhalt  durch Gustav Schwarzkopf am 11. 8. 1899 in Tiefer Graben 23}\toendnotes[C]{\smallbreak}
\Standort{DLA, A:Schnitzler, HS.1985.1.1897.}
\physDesc{Bildpostkarte, 240 Zeichen
\newline{}Handschrift Arthur Schnitzler: Bleistift, deutsche Kurrent
\newline{}Handschrift Richard Beer-Hofmann: Bleistift, lateinische Kurrent
\newline{}Versand: 1) Stempel: »\nobreak{}\oindex{Predazzo@\textbf{Predazzo}|pwk}Pre\textcolor{gray}{d}azzo, 9. 8. 99\nobreak{}«.   2) Stempel: »\nobreak{}\oindex{I., Innere Stadt@\textbf{I., Innere Stadt}|pwk}Wien 1/1 1, 11. 8. 1899, 9–10½V., Bestellt\nobreak{}«. }\pstart{}{\pb}Herrn \textsc{Gustav Schwarzkopf}\pend{}\pstart{}Wien\oindex{Wien@\textbf{Wien}|pw}\pend{}\pstart{}\textsc{Tiefer Graben 23}\oindex{Wien@\textbf{Wien}!I., Innere Stadt@\textbf{I., Innere Stadt}!Tiefer Graben 23@\textbf{Tiefer Graben 23}|pw}\pend{}{\bigskip}
\pstart
           \centering{}{\pb}\textcolor{gray}{\textbf{Un Saluto da PREDAZZO\oindex{Predazzo@\textbf{Predazzo}|pw}.}}\pend
           
\pstart
           \centering{}\textcolor{gray}{\textbf{Predazzo\oindex{Predazzo@\textbf{Predazzo}|pw}
                  m. Lattemar\oindex{Latemar@\textbf{Latemar}|pw}.}}\pend
           \vspace{1em}
\pstart
           {\pb}Mittwoch 9/8 99.\pend
           \vspace{0.5em}
\pstart
           Montag über Giau\oindex{Passo di Giau@\textbf{Passo di Giau}|pw}, und \textsc{Caprile\oindex{Caprile@\textbf{Caprile}|pw}}{ }geſtern über \textsc{Fedaja\oindex{Fedaia@\textbf{Fedaia}|pw}} nach \textsc{Vigo\oindex{Vigo di Fassa@\textbf{Vigo di Fassa}|pw}} –
      alles zu Fuſs!! Heute nicht
      wegen Regen, mit der Post hieher.\pend
           
\pstart
           Viele herzliche Grüße.{\\[\baselineskip]}\spacefill\mbox{Arthur}\pend
           \leftskip=0em{}\selectlanguage{ngerman}\vspace{1em}\pstart \spacefill\mbox{{[}hs. Beer-Hofmann:{]} Richard}\pend{}
\pstart
           \noindent{}Wassermann\pwindex{Wassermann, Jakob 10.\,3.\,1873 Fürth – 1.\,1.\,1934 Altaussee@\textsc{Wassermann, Jakob} (10.\,3.\,1873 Fürth – 1.\,1.\,1934 Altaussee)|pw} grüsst – aber wirklich!\pend
           \selectlanguage{ngerman}\endnumbering\briefempfaengerindex{Schwarzkopf, Gustav@\textsc{Schwarzkopf, Gustav}!zzzBeer-Hofmann, Richard@\emph{von Richard Beer-Hofmann}!1899-08-091@{9. 8. 1899}|)be}\briefempfaengerindex{Schwarzkopf, Gustav@\textsc{Schwarzkopf, Gustav}!zzzSchnitzler, Arthur@\emph{von Arthur Schnitzler}!1899-08-091@{9. 8. 1899}|)be}\mylabel{L04435h}
\begin{anhang}
\end{anhang}\newcommand{\dateiname}{L04435}\newcommand{\titel}{Arthur Schnitzler und Richard Beer-Hofmann an Gustav Schwarzkopf, 9. 8. 1899}\newcommand{\editorInnen}{Herausgegeben von Jahnke, SelmaMüller, Martin Anton}\input{../tex-inputs/latex-leseansicht-abspann}
